\section*{Введение}
\addcontentsline{toc}{section}{Введение}

Теория вероятностей представляет собой фундаментальный раздел математики, изучающий закономерности не просто в отдельных случайных явлениях, а именно в \textbf{массовых случайных явлениях}. Только при наличии достаточно большой массы наблюдений, проводимых в одних и тех же условиях, становится возможным установление устойчивых закономерностей в случайных процессах. 

Рассмотрим классический пример: бросание монеты. Выпадение «герба» или «решки» в результате единичного броска является событием чисто случайным. Теория вероятностей не ставит своей целью объяснить, почему в данном конкретном случае выпал «герб», а не «решка». Ее фундаментальная задача заключается в ином: определить, насколько вероятно выпадение «герба» при условии, что монета будет подброшена достаточно большое количество раз. Именно в пределе большого числа испытаний хаотичность отдельных событий уступает место строгим статистическим закономерностям.

К созданию и развитию теории вероятностей и математической статистики приложили свои усилия многие выдающиеся ученые. Среди пионеров и классиков этой науки следует назвать Г.~Галилея, Б.~Паскаля, П.~Ферма, Х.~Гюйгенса, Я.~Бернулли, А.~Муавра, П.~Лапласа, С.~Пуассона, а также ученых, внесших вклад в ее современное развитие: У.~Госсета (Стьюдента), К.~Пирсона, Н.~Винера, У.~Феллера, Дж.~Дуба, Р.~Фишера и многих других.

Особый вклад в становление и строгое аксиоматическое обоснование теории вероятностей внесли наши соотечественники. Славная плеяда российских и советских математиков включает В.~Я.~Буняковского, П.~Л.~Чебышева, А.~М.~Ляпунова, А.~А.~Маркова, С.~Н.~Бернштейна, А.~Я.~Хинчина, А.~Н.~Колмогорова, В.~И.~Романовского, Н.~В.~Смирнова и многих других исследователей, чьи работы заложили незыблемый фундамент современной вероятностно-статистической науки.

\vspace{1em}
\noindent\textit{Данная методичка призвана помочь вам систематизировать знания по этим темам, освоить ключевые методы анализа распределений и научиться применять их для решения практических задач.}